\clearpage
\phantomsection% 
\addcontentsline{toc}{chapter}{ABSTRAK}
\thispagestyle{fancy}

\begin{center}
    \large \bfseries \MakeUppercase{Abstrak}\\
    \normalsize \normalfont {\thetitle}\\
    \normalsize \normalfont {\theauthor}\\
    \bigskip
    
    \normalsize \normalfont \justifying \singlespacing
    Prevalensi \textit{stunting} di Indonesia merupakan masalah kesehatan komunal yang dipengaruhi oleh berbagai faktor determinan yang kompleks dan non-linear. Penelitian ini bertujuan untuk mengidentifikasi faktor penentu utama \textit{stunting} tingkat provinsi dan mengekstraksi titik ambang batas kritisnya menggunakan agregasi data Survei Status Gizi Indonesia (SSGI) 2024. Metodologi pemodelan dibangun menggunakan algoritma \textit{Random Forest Regression}, sementara interpretabilitas model untuk membuka sifat \textit{black-box} dianalisis melalui pendekatan \textit{Explainable AI} menggunakan metode \textit{SHapley Additive exPlanations} (SHAP). Hasil evaluasi menunjukkan model memiliki kinerja prediktif yang sangat baik, dibuktikan dengan perolehan nilai koefisien determinasi ($R^2$) sebesar 0,9090, \textit{Mean Absolute Error} (MAE) sebesar 1,40, \textit{Root Mean Squared Error} (RMSE) sebesar 1,77, dan tingkat kesalahan prediksi (MAPE) sebesar 7,22\%. Berdasarkan visualisasi tingkat kepentingan fitur SHAP, faktor determinan yang paling mendominasi adalah kurangnya pemenuhan standar gizi harian balita berupa ketiadaan \textit{Minimum Dietary Diversity} (MDD) dan \textit{Minimum Milk Feeding Frequency} (MMFF), tingginya persentase usia ibu berisiko, akses infrastruktur (sanitasi berisiko), serta minimnya Pemberian Makanan Tambahan (PMT). Lebih lanjut, analisis \textit{Dependence Plot} berhasil mengungkap ambang batas toleransi matematis komunal, seperti ketiadaan MDD di angka 55\%, ketiadaan MMFF 22\%, usia ibu berisiko 14\%, dan fasilitas sanitasi berisiko pada tingkat 10\%. Lonjakan risiko prevalensi \textit{stunting} terjadi secara drastis apabila suatu wilayah melampaui batas persentase tersebut. Kesimpulannya, model prediksi dan penjelasan visual yang dihasilkan terbukti efektif memetakan akar masalah malnutrisi secara presisi, sehingga dapat dimanfaatkan sebagai acuan dasar evaluasi kebijakan penanganan \textit{stunting} nasional.\par

    \vspace{20pt}
    \raggedright \textbf{Kata Kunci: \textit{Stunting}, \textit{Random Forest}, SHAP, SSGI, Determinan}
    
    \vfill
    
\end{center}
\clearpage